%% kapa: exam=1 round=19 year=2008 period=2 form=A subject=economics subject_ko=경제원론 seq=32 number=32
%% source: Q-Net 감정평가사 기출문제 (cst00309) artlSeq=1001962 file=extracted/1008/1001962/2008년 제19회 감정평가사 1차시험 기출문제/2교시.hwp
%% conversion: kordoc 4.24.51 -> tools/gen1.py (auto); review=proofread-sonnet; flags=none
\begin{question}{32}
다음 중 화폐의 중립성에 대한 설명으로 옳은 것을 모두 고르면?
\begin{framedtext}
ㄱ. 케인즈학파는 화폐의 중립성이 단기에 성립한다고 주장한다. \par
ㄴ. 고전학파는 화폐의 중립성이 장기에서는 물론 단기에서도 성립한다고 주장한다. \par
ㄷ. 통화공급이 변화할 때 실질변수가 영향을 받지 않는 경우를 말한다. \par
ㄹ. 임금을 비롯한 가격변수가 경직적인 경우에 성립하지 않는다. \par
\end{framedtext}
\begin{choices}
  \item ㄱ, ㄷ
  \item ㄴ, ㄹ
  \item ㄷ, ㄹ
  \item ㄱ, ㄴ, ㄷ
  \item ㄴ, ㄷ, ㄹ
\end{choices}
\end{question}
\answer{5}
